\begin{table}[H]
\centering
\caption{The Depression Shield: Migration and Occupational Resilience, 1930--1940}
\label{tab:main}
\begin{threeparttable}
\begin{tabular}{lccc}
\toprule
 & (1) & (2) & (3) \\
 & OLS & IV & IV \\
\midrule
\multicolumn{4}{l}{\textit{Panel A: Second stage --- Dep.\ var: $\Delta$ Occ.\ Score, 1930--1940}} \\
\addlinespace
  Migrant (S$\rightarrow$N) & -0.214^{***} & 2.512^{**} & 3.167^{***} \\
   & (0.061) & (0.980) & (0.911) \\
\addlinespace
\midrule
\multicolumn{4}{l}{\textit{Panel B: First stage --- Dep.\ var: Migrant (S$\rightarrow$N)}} \\
\addlinespace
  Log(Distance$^{-1} \times$ Black Pop.\ 1910) & & 0.0452^{***} & 0.0452^{***} \\
   & & (0.0039) & (0.0037) \\
  First-stage $F$ & & 137.0 & 152.3 \\
\addlinespace
\midrule
  Individual controls & $\checkmark$ & & $\checkmark$ \\
  Observations & 222,380 & 222,380 & 222,380 \\
\bottomrule
\end{tabular}
\begin{tablenotes}[flushleft]
\small
\item \textit{Notes:} Sample is Black males aged 15--55 in 1920, born in Southern states, linked across three Census decades via the IPUMS MLP. Migrant equals one if the individual resided in the South in 1920 and a Northern/Western state by 1930. The instrument is the log of a shift-share measure combining inverse geographic distance from the origin county to 12 major Northern cities with each city's 1910 Black population stock. This measure captures pre-existing migration corridors: counties closer to Northern cities with established Black communities had higher out-migration rates. Individual controls: age, age squared, school attendance, farm worker status, marital status, and baseline occupational score (all measured in 1920). Standard errors clustered at origin county in parentheses. $^{***}$p$<$0.01, $^{**}$p$<$0.05, $^{*}$p$<$0.10.
\end{tablenotes}
\end{threeparttable}
\end{table}

